%% Flagship, full-length. Target: arXiv (cs.SE, cross-list cs.PL), NAMED.
%% `nier/main.tex` is a 4+1 page anonymised cut of this; `demo/main.tex` is the
%% separate tool paper. No venue page limit here — build.mjs reports, never fails.
%%
%% ARGUMENT: paper/flagship/OUTLINE.md is the authority. Read it before editing.
%% CLAIM SAFETY: paper/README.md lists what may not be said. Read that too.
\documentclass[10pt,journal,compsoc]{IEEEtran}

\usepackage{cite}
\usepackage{amsmath}
\usepackage{graphicx}
\usepackage{booktabs}
\usepackage{xspace}
\usepackage{url}
\usepackage[hidelinks]{hyperref}
\usepackage{listings}
\usepackage{xcolor}

%% Structural counts are INTERPOLATED from the compiler's own tables by
%% `npm run gen:paper-facts` and gated by `npm run check:drift`. Never retype a
%% number into this file: doing so is the defect this paper is about, and the
%% project's own prose made that mistake three times before this file existed.
% GENERATED by scripts/gen-paper-facts.ts — do not hand-edit.
% Every number here is derived from the compiler's own source of truth.
% Regenerate with `npm run gen:paper-facts`; `npm run check:drift` gates it.
\newcommand{\factCliCommands}{20\xspace}
\newcommand{\factCliExamples}{46\xspace}
\newcommand{\factCliFlags}{115\xspace}
\newcommand{\factDoorKinds}{5\xspace}
\newcommand{\factErrorCodesError}{73\xspace}
\newcommand{\factErrorCodesTotal}{115\xspace}
\newcommand{\factErrorCodesWarning}{42\xspace}
\newcommand{\factExportFormats}{5\xspace}
\newcommand{\factExportFormatsZeroDep}{3\xspace}
\newcommand{\factKeywordsAttribute}{54\xspace}
\newcommand{\factKeywordsControl}{20\xspace}
\newcommand{\factKeywordsElement}{11\xspace}
\newcommand{\factKeywordsEnum}{51\xspace}
\newcommand{\factKeywordsTotal}{136\xspace}
\newcommand{\factLintRules}{21\xspace}
\newcommand{\factOperators}{21\xspace}

%% HAND-WRITTEN, and the only file in this directory where that is allowed.
%%
%% These are external or paid measurements that cannot be derived from the
%% repository, so `gen:paper-facts` cannot produce them and `check:drift` cannot
%% gate them. The discipline that replaces the gate: EVERY entry carries a comment
%% naming its source and the date it was measured. An entry without provenance does
%% not go in this file.
%%
%% Everything structural belongs in facts.tex (generated, drift-gated); everything
%% about repository scale belongs in scale-snapshot.json (generated on demand with
%% an explicit --date).

% ---------------------------------------------------------------------------
% Repository scale. Source: paper/scale-snapshot.json, commit 8096f2d, 2026-08-22.
% Regenerate with: npm run paper:snapshot -- --date YYYY-MM-DD
% ---------------------------------------------------------------------------
\newcommand{\factSnapshotDate}{2026-08-22\xspace}
\newcommand{\factSnapshotCommit}{8096f2d\xspace}
\newcommand{\factVersion}{1.26.1\xspace}
\newcommand{\factCommits}{513\xspace}
\newcommand{\factTags}{43\xspace}
\newcommand{\factFirstCommit}{2026-06-25\xspace}
\newcommand{\factDevDays}{52\xspace}          % 2026-06-25 .. 2026-08-16
\newcommand{\factRepoLines}{180{,}749\xspace}
\newcommand{\factSrcLines}{36{,}207\xspace}
\newcommand{\factTestLines}{29{,}204\xspace}
\newcommand{\factTestFiles}{163\xspace}
\newcommand{\factExamples}{27\xspace}
\newcommand{\factAdrs}{16\xspace}

% ---------------------------------------------------------------------------
% GBNF grammar/parser agreement.
% Source: paper/experiments/gbnf/agreement-results.json + NUMBERS-AUDIT.md Sec.2.
% Reproduce: npx tsx paper/experiments/gbnf/agreement.mts --rev 2ac49ef^ --rev v1.25.0
% NOTE: the project's own prose said "71-plan corpus" (wrong; it was never 71) and
% "113/113" (a whole-FILE it() count that grows with the example gallery, today 126).
% Cite only the figures below.
% ---------------------------------------------------------------------------
\newcommand{\factGbnfCorpus}{73\xspace}
\newcommand{\factGbnfParseable}{36\xspace}
\newcommand{\factGbnfRejected}{37\xspace}
\newcommand{\factGbnfAgreeNow}{73\xspace}
\newcommand{\factGbnfAgreeBefore}{49\xspace}
\newcommand{\factGbnfDisagree}{24\xspace}
\newcommand{\factGbnfOverDerive}{21\xspace}    % grammar derived, parser rejects
\newcommand{\factGbnfUnderDerive}{3\xspace}    % parser accepts, grammar cannot derive

% ---------------------------------------------------------------------------
% Mutation experiment (NEW for this paper — not the project's earlier,
% unreproducible "27 of 28", which is deliberately not cited anywhere).
% Source: paper/experiments/mutation/results.json, window 01:11:13Z-01:16:18Z on
% 2026-08-22, node v24.15.0, vitest ^2.1.8, commit 8096f2d.
% Reproduce: node paper/experiments/mutation/run.mjs
%            npx tsx paper/experiments/mutation/witness.mts
%
% PROVENANCE, and it is a finding in its own right. An EARLIER run of this same
% harness over the same subject reported 49/56 (87.5%). It is VOID: its window was
% 00:41:16Z-00:46:24Z and a third party restored one of the mutated files from HEAD
% at ~00:44, mid-flight, so two mutants were counted as survivors purely because
% their mutation had been removed underneath them. The artifact recorded
% `treeCleanAfterRun: true` over that run. Cite ONLY the figures below; the void
% run appears in the paper as a worked example of a corrupted measurement producing
% a plausible-but-wrong number, never as a result.
% ---------------------------------------------------------------------------
\newcommand{\factMutants}{56\xspace}
\newcommand{\factMutantsMeasure}{20\xspace}
\newcommand{\factMutantsFrame}{36\xspace}
\newcommand{\factMutantsKilled}{51\xspace}
\newcommand{\factMutantsSurvived}{5\xspace}
\newcommand{\factMutantScore}{91.1\%\xspace}
\newcommand{\factSurvivorsEquivalent}{5\xspace}   % all five; the two real gaps are now killed
\newcommand{\factMutantsVoidKilled}{49\xspace}    % the corrupted run — example only, never a result
\newcommand{\factMutantScoreVoid}{87.5\%\xspace}

% ---------------------------------------------------------------------------
% Property-generator repair. Source: CHANGELOG.md v1.26.1 / AGENTS.md.
% The old fc.string() body produced zero geometry across 5000 samples; the
% replacement kills a planted lowerWalls iteration-order bug in 5 runs that the old
% property survived 300 runs of.
% ---------------------------------------------------------------------------
\newcommand{\factOldPropSamples}{5000\xspace}
\newcommand{\factOldPropGeometry}{0\xspace}
\newcommand{\factNewPropSamples}{3000\xspace}
\newcommand{\factPlantSurvivedRuns}{300\xspace}
\newcommand{\factPlantKilledRuns}{5\xspace}
\newcommand{\factWatchBrokenReleases}{25\xspace}   % arch watch, v1.1.0 -> v1.26.0

% ---------------------------------------------------------------------------
% Authorability baseline. CANONICAL RUN = eval/live-baseline.json.
% OpenAI gpt-5.5-2026-04-23, seed 20260711 (requested, NOT a determinism guarantee),
% max_completion_tokens 16384, 26 briefs, judge v2 / synonyms v1,
% GitHub Actions run 29190294073, 2026-07-12.
% The per-brief scorecard and Gate G1's control arm belong to the SUPERSEDED
% 2026-07-11 run (29150982395) — see facts below, and never mix the two.
% ---------------------------------------------------------------------------
\newcommand{\factEvalBriefs}{26\xspace}
\newcommand{\factEvalValid}{23\xspace}
\newcommand{\factEvalIntent}{14\xspace}
\newcommand{\factEvalSound}{3\xspace}
\newcommand{\factEvalIntentLone}{18\xspace}    % L1 deterministic overlay, zero extra model calls
\newcommand{\factEvalSoundLone}{7\xspace}
\newcommand{\factEvalRepairMoves}{41\xspace}
% The judge recalibration: v1 -> v2 moved intent 9% -> 50% with ZERO model change.
% Pre-registered prediction was 55-65% artifact / true rate 45-60%; outcome 50-54%.
\newcommand{\factJudgeBefore}{9\%\xspace}
\newcommand{\factJudgeAfter}{50\%\xspace}

% ---------------------------------------------------------------------------
% Gate G1 (NL brief -> intent-JSON faithfulness). Source: eval/g1/report.md,
% scores-{model,fable,human}.json; all figures re-verified 2026-08-21.
% PROSE REQUIREMENT: this was graded by TWO INDEPENDENT MODEL RATERS, with a human
% adjudicating only the two disagreements. No human graded any assertion. The
% sensitivity result MUST travel in the same breath as the headline z.
% ---------------------------------------------------------------------------
\newcommand{\factGoneFaithful}{154\xspace}
\newcommand{\factGoneAssertions}{157\xspace}
\newcommand{\factGoneRate}{98.1\%\xspace}
\newcommand{\factGoneControl}{93.4\%\xspace}       % 155/166, deliverable semantics
\newcommand{\factGoneZ}{2.08\xspace}
\newcommand{\factGoneP}{0.019\xspace}
\newcommand{\factGoneControlValidOnly}{95.7\%\xspace}  % 155/162
\newcommand{\factGoneZValidOnly}{1.24\xspace}
\newcommand{\factGonePValidOnly}{0.11\xspace}
\newcommand{\factGoneKappa}{0.50\xspace}

% ---------------------------------------------------------------------------
% Distribution availability. Source: paper/experiments/npm-downloads/, npm registry
% API, fetched 2026-08-21 (range ends 2026-08-20, the last complete day).
% USAGE RULE: the release-cadence concentration must appear in the SAME SENTENCE as
% the download figure. All five of the core package's busiest days are days a v* tag
% was pushed, and they carry 44% of the 90-day total. Never call these users,
% installs or adoption.
% ---------------------------------------------------------------------------
\newcommand{\factNpmThirty}{2{,}357\xspace}
\newcommand{\factNpmNinety}{8{,}153\xspace}
\newcommand{\factNpmMcpThirty}{1{,}181\xspace}
\newcommand{\factNpmTopFiveShare}{44\%\xspace}
\newcommand{\factNpmFirstDownload}{2026-06-24\xspace}
\newcommand{\factNpmFetched}{2026-08-21\xspace}

% ---------------------------------------------------------------------------
% ECOSYSTEM CENSUS — the paper's principal empirical result.
% Source: paper/experiments/ecosystem/results-2026-08-22.json + FINDINGS.md.
% Registry enumerated 2026-08-22T00:36:08Z, 242 pages. A CENSUS of registry-listed
% MCP servers with an npm package, not a sample — so no sampling error in M1.
% Reproduce/verify: node paper/experiments/ecosystem/verify-findings.mjs
% USAGE RULES: the headline is the CORROBORATED figure; the raw figure is the upper
% bound and must be labelled as such. `never-published` is NOT corroborable by
% construction and is excluded from the headline, so the headline understates.
% Never quote the median without the tail, or the tail without the median.
% ---------------------------------------------------------------------------
\newcommand{\factEcoServers}{24{,}131\xspace}
\newcommand{\factEcoPackages}{7{,}272\xspace}
\newcommand{\factEcoApplicable}{5{,}052\xspace}
\newcommand{\factEcoDerived}{525\xspace}
\newcommand{\factEcoDerivedPct}{10.4\%\xspace}
\newcommand{\factEcoRetyped}{4{,}527\xspace}
\newcommand{\factEcoRetypedPct}{89.6\%\xspace}
\newcommand{\factEcoDriftRaw}{2{,}777\xspace}
\newcommand{\factEcoDriftRawPct}{55.0\%\xspace}       % upper bound
\newcommand{\factEcoDriftConfirmed}{2{,}075\xspace}
\newcommand{\factEcoDriftConfirmedPct}{41.1\%\xspace} % headline
\newcommand{\factEcoIndependentVersioning}{68\xspace} % correctly excluded — the check working
\newcommand{\factEcoNeverPublished}{578\xspace}
% Of the never-published values, the share matching a common scaffold default
% (1.0.0, 0.1.0, 0.0.1, 0.0.0, 0.1.1, 0.0.2, 1.0.1, 1.0.0-beta). `1.0.0` alone is
% 39.4% and `0.1.0` is 22.8%. Computed from results-2026-08-22.json, 2026-08-22.
% This licenses "typically a scaffold default" and NOTHING about when it was
% written — we did not inspect commit history and must not imply that we did.
\newcommand{\factEcoNeverPublishedScaffoldPct}{71.6\%\xspace}
\newcommand{\factEcoMedianReleases}{2\xspace}
\newcommand{\factEcoMaxReleases}{260\xspace}
\newcommand{\factEcoMaxDays}{497\xspace}
\newcommand{\factEcoGeFive}{337\xspace}
\newcommand{\factEcoGeTen}{131\xspace}
\newcommand{\factEcoInternalDisagreement}{112\xspace}
\newcommand{\factEcoFpSample}{30\xspace}
\newcommand{\factEcoFpFound}{0\xspace}
\newcommand{\factEcoFpUpperBound}{10\%\xspace}        % rule of three, n=30, 0 events
% Confidence intervals. These are binomial intervals on the underlying propensity,
% NOT sampling error — M1 is a census. Say so whenever one is quoted, or a reviewer
% will reasonably ask why a census has a CI.
\newcommand{\factEcoDriftConfirmedCI}{39.7--42.4\%\xspace}
\newcommand{\factEcoDriftRawCI}{53.6--56.3\%\xspace}
% Pooled false-positive review over the CORROBORATED stratum (the headline's own
% population), including the 30 hand-reviewed independently. False positives
% concentrate in `never-published` and Tier B, neither of which feeds the headline.
\newcommand{\factEcoFpReviewed}{130\xspace}
\newcommand{\factEcoFpRate}{0.8\%\xspace}
\newcommand{\factEcoFpRateCI}{0.1--4.2\%\xspace}
% Metric 2 — a SAMPLE (600 of 6,403 eligible), reported separately, never merged into M1.
\newcommand{\factEcoMTwoCompared}{267\xspace}
\newcommand{\factEcoMTwoDiffer}{39\xspace}
\newcommand{\factEcoMTwoDifferPct}{14.6\%\xspace}

% ---------------------------------------------------------------------------
% Downstream impact of the drifted specification. Source: AGENTS.md v1.26.0 entry,
% measured in the ArchCanvas consumer.
% ---------------------------------------------------------------------------
\newcommand{\factSpecBadLines}{7\xspace}
\newcommand{\factSpecUncompilable}{4\xspace}
\newcommand{\factDownstreamFailed}{11\xspace}
\newcommand{\factDownstreamTotal}{18\xspace}
\newcommand{\factGbnfIllegalForms}{11\xspace}   % forms the shipped grammar derived that the parser rejects

% ---------------------------------------------------------------------------
% Executable self-description (the remedy, as shipped). Source: test/spec-forms.test.ts.
% ---------------------------------------------------------------------------
\newcommand{\factSpecFormsCompile}{44\xspace}
\newcommand{\factSpecFormsRefused}{19\xspace}

% ---------------------------------------------------------------------------
% llms.txt ecosystem context. Source: Originality.AI llms.txt tracking study and
% ppc.land coverage, accessed 2026-08-22. Order-of-magnitude only — cite as "more
% than 800,000 sites", never to three significant figures.
% ---------------------------------------------------------------------------
\newcommand{\factLlmsTxtSites}{800{,}000\xspace}


\lstdefinelanguage{archlang}{
  morekeywords={plan,wall,room,door,window,opening,furniture,dim,level,zone,place,site,paper,axes,schedule,legend,component,import},
  morecomment=[l]{//},
  morestring=[b]",
}
\lstset{
  language=archlang, basicstyle=\ttfamily\footnotesize,
  keywordstyle=\color{blue!60!black}, commentstyle=\color{black!45},
  columns=fullflexible, keepspaces=true, showstringspaces=false,
  aboveskip=4pt, belowskip=4pt,
}

\begin{document}

\title{Reproducible Is Not Correct: Reflexive Gates and the\\Self-Descriptions We Ship to Language Models}

\author{Chan~Meng%
\IEEEcompsocitemizethanks{\IEEEcompsocthanksitem C. Meng is an independent researcher.
E-mail: ai@gavigo.com}}

\markboth{Preprint}{Meng: Reproducible Is Not Correct}
\IEEEtitleabstractindextext{%
\begin{abstract}
Software increasingly ships descriptions of itself intended for a language model
rather than a person: specifications, constrained-decoding grammars, JSON schemas,
vendored context files, and the identity a server states to its host at handshake
time. These artifacts are neither compiled nor executed against the thing they
describe. The gate deployed over them, almost universally, is \emph{reflexive}:
regenerate the artifact and fail on a difference from the committed copy. We observe
that such a gate establishes determinism and cannot, even in principle, establish
truth --- and that it is nonetheless sound wherever it is deployed, because the
artifacts it guards are \emph{mechanically derived} from their subject, and a derived
artifact can be stale but never wrong. That precondition is never stated and never
checked. When a generator \emph{retypes} a fact instead of deriving it, the same gate
is vacuous, and it looks identical in continuous integration.
We measure what this costs. A census of all \factEcoServers{} servers listed on the
Model Context Protocol registry, covering the \factEcoPackages{} unique npm packages
among them, finds that of the \factEcoApplicable{} packages whose self-reported
identity can be resolved statically, \factEcoRetypedPct{} retype it rather than derive
it, and \factEcoDriftConfirmedPct{} state a version that once tracked the package and
has since been left behind --- corroborated against each package's own release
history, with independently versioned cases excluded. We add six instances across
four independent third-party compilers, four of them verified by execution, and a
depth study of one production language whose four self-descriptions were all wrong on
first check. We claim none of the underlying concepts, all of which are established;
we report the precondition, the artifact class, and the measurement. The subject
system was substantially co-authored with language models, and we discuss what that
implies for which guards matter.
\end{abstract}
\begin{IEEEkeywords}
Software verification, test adequacy, documentation consistency, build systems,
domain-specific languages, LLM code generation.
\end{IEEEkeywords}}
\maketitle
\IEEEdisplaynontitleabstractindextext
\IEEEpeerreviewmaketitle

\section{Introduction}\label{sec:intro}

\IEEEPARstart{A}{} growing share of what software projects ship is a description of
themselves addressed to a machine. A package that wants to be usable by a language
model publishes a specification for the model to read, a grammar for a constrained
decoder to sample under, JSON schemas for its inputs, context files bundled at pack
time, and --- if it speaks a protocol --- an identity it states at handshake time.
None of this is read by a person in the ordinary course of events. None of it is
executed. Nothing runs it against the system it purports to describe.

What guards it, when anything does, is a \emph{drift gate}: regenerate the artifact
from its generator and fail the build on any difference from the committed copy.
This pattern is close to universal wherever code generation happens. Microsoft's API
Extractor regenerates a public-API report and fails a production build on any
difference from the checked-in file~\cite{api_extractor}; Bazel's Gazelle offers
\texttt{-mode=diff} for exactly this~\cite{bazel_gazelle}; Kubernetes projects
regenerate CRD manifests and diff them; \texttt{terraform fmt -check} and
\texttt{cargo readme -{}-check} are the same shape. The .NET source-generator
documentation recommends snapshot testing against a committed snapshot and justifies
it explicitly on the grounds that generator output is deterministic for a given
input.

That justification is the whole subject of this paper. A reflexive gate --- one that
compares a generator to its own output --- establishes that the generator is
deterministic. It cannot establish that what the generator emits is \emph{true}. Both
halves of this are already in the literature, stated separately and never joined:
build-system theory places the correctness of generated content outside what a build
system guarantees~\cite{mokhov2020buildsystems}, and the reproducible-builds programme
adopts byte-reproducibility explicitly as an \emph{integrity} property rather than a
correctness one~\cite{lamb2022reproducible}.

The reason the gap has not needed naming is that, until recently, it did not matter.
Every reflexive gate in the list above guards an artifact that is \emph{mechanically
derived} from its subject: API Extractor reads the TypeScript compiler's own symbol
table, Gazelle reads real Go imports, \texttt{protoc} reads the \texttt{.proto},
\texttt{cargo-semver-checks} reads rustdoc JSON. \textbf{A derived artifact cannot be
wrong about its subject --- only stale.} For derived artifacts the reflexive gate is
therefore sufficient, and its limitation is invisible because it has no consequences.

Our observation is that this soundness rests on a precondition that is never stated
and never checked, and that the LLM-facing artifact class routinely violates it. A
generator that \emph{retypes} a language fact into a specification, or hand-writes a
grammar's productions beside the parser rather than deriving them from it, produces an
artifact that can be confidently, stably, reproducibly wrong. It wears the same gate.
It looks identical in CI. And unlike a stale comment, which misleads a reader who can
notice, a wrong machine-consumed description makes invalid output possible \emph{by
construction}: a constrained-decoding grammar exists for the single purpose of making
syntactically invalid output impossible, and one that over-derives has inverted its own
function while continuing to appear to perform it.

\subsection{What we do not claim}\label{sec:concede}

Almost every component of what follows is someone else's, and we concede it here
rather than have it supplied.

A guard that passes without its passing depending on what it appears to check is
\emph{vacuous satisfaction}, carried from model checking into testing by Ball and
Kupferman~\cite{ball2008vacuitytesting}; the remedy we use --- perturb the thing under
test and require the guard to notice --- is Beer et al.'s vacuity detection procedure,
published in 1997~\cite{beer1997vacuity}. Executing documentation against the
implementation it describes is older than it looks: \texttt{@tComment} turned Javadoc
into runnable assertions in 2012~\cite{tan2012tcomment}, CASCADE now generates tests
directly from documentation~\cite{kiecker2026cascade}, and Rust, Go and Python all run
their doc examples by default~\cite{rust_doctests,go_testable_examples,python_doctest}.
Requiring that a tool's declared diagnostics can actually fire is shipped and
CI-enforced: rustc's \texttt{tidy} check demands that every error code be documented,
triggered by a test, and emitted by the compiler~\cite{rustc_tidy_errorcodes}.
Validating a machine-consumed \emph{description} against the system it describes is a
decade-old industry practice --- Pact generates a contract during consumer tests and
verifies it by executing it against the real provider~\cite{pact_contract_testing}, and
Schemathesis runs a schema against a live server and reports divergence without
deciding which side is wrong~\cite{schemathesis}. Even the remedy we converge on has a
name: deriving a fact once instead of retyping it is DRY~\cite{hunt1999pragmatic}.

What is left is a precondition and a measurement. DRY says what you should do; nobody
appears to have stated the failure mode of the gate when DRY is violated --- that
regenerate-and-diff is vacuous for a retyped generator, that it cannot detect the
violation, and that it looks identical in CI. DRY is a design rule with no enforcement
story; the reflexive gate is an enforcement mechanism whose soundness silently depends
on that rule and never checks it. We report what the junction costs, in the one
artifact class with no executable counterpart to fall back on, and we report it as
measurement, not as definition.

\subsection{Contributions}

\begin{itemize}
\item The \emph{derived-versus-retyped} precondition of the reflexive gate, with an
  empirical demonstration that the gate cannot see the difference
  (Section~\ref{sec:mechanism}).
\item An ecosystem census of \factEcoPackages{} packages measuring how often a
  machine-consumed self-description is false, with a corroboration step that separates
  drift from deliberate independent versioning (Section~\ref{sec:census}).
\item Six instances across four independent third-party compilers, four verified by
  execution rather than by reading, together with four refutations
  (Section~\ref{sec:corroboration}).
\item A depth study of one production language, including two measurements that were
  themselves corrupted and what that cost (Section~\ref{sec:depth}).
\item A remedy stated as a requirement --- every guard must be provably non-vacuous ---
  and applied to this paper's own toolchain, which produced six vacuous guards while
  the paper was being written (Sections~\ref{sec:remedy} and~\ref{sec:self}).
\end{itemize}

\section{The Mechanism}\label{sec:mechanism}

\subsection{Two kinds of generator}

Let $S$ be a system, $G$ a generator, and $A = G(S)$ a committed artifact describing
$S$. A reflexive gate recomputes $G(S)$ at build time and fails if it differs from the
committed $A$. Write $A^{\ast}$ for the artifact that would be true of $S$.

Call $G$ \emph{derived} when its output is a function of $S$ alone --- when it reads
the subject and reports what it finds. For a derived generator, $G(S) = A^{\ast}$ holds
by construction: the artifact cannot disagree with the subject, because the subject is
the only thing it consulted. The single failure available to it is \emph{staleness} ---
$A$ was correct for an older $S$ and has not been recomputed --- and that is exactly
what a reflexive gate detects, because recomputing $G$ against the current $S$ produces
a different result.

Call $G$ \emph{retyped} when its output depends on facts a human transcribed into the
generator. Now $G(S)$ and $A^{\ast}$ are independent: the generator emits what its
author believed, deterministically, forever. Regenerating produces the same bytes.
The gate passes. The artifact is wrong.

The distinction is not exotic. It is the difference between a generator that walks the
compiler's own keyword table and one whose template contains a hand-written list of
keywords that a maintainer keeps up to date. Both are generators, both are covered by
the same \texttt{npm run gen \&\& git diff -{}-exit-code} step, both are green.

\subsection{Why the gate cannot check its own precondition}

A reflexive gate observes exactly one thing: whether $G(S)$ equals the committed $A$.
That comparison is blind to how $G$ computed its answer. Two generators over the same
subject, one derived and one retyped, produce identical CI output --- a passing step
with no diff --- and no signal distinguishes them. The gate has no access to the
property on which its own soundness depends.

Nor is the precondition usually apparent to a reader. The generators in the subject
system of Section~\ref{sec:depth} are several hundred lines each and mix both styles:
they derive the keyword vocabulary from a shared table \emph{and} carry hand-written
prose that states grammar rules in English. A single generator can be derived in one
paragraph and retyped in the next.

\subsection{The demonstration}\label{sec:twogenerators}

The cleanest evidence that the gate is blind is a case where two generators disagreed
with the parser \emph{and with each other}, in different places, while both gates
stayed green for three releases.

The subject system generates two artifacts that describe its language: a one-page
specification for a model to read, and a GBNF grammar for a constrained decoder to
sample under. Both are produced by generators; both are covered by the drift gate; both
had been green since they were written. When the two artifacts were first checked
against the parser rather than against themselves, the specification taught
\factSpecBadLines{} incorrect grammar lines of which \factSpecUncompilable{} do not
compile, and the grammar derived \factGbnfIllegalForms{} forms the parser rejects. The
sets did not coincide. Each generator was right where the other was wrong.

Two observations follow. First, no amount of re-running either gate could have found
this: both generators reproduced their own output perfectly, which is all the gate ever
asked. Second, the failure was not detectable by comparing the artifacts to each other
either --- they disagreed, but with no oracle there is no way to say which was right,
and in fact both were wrong in different places. The only thing that settles it is
executing them against the parser.

The downstream cost was measurable. A separate consumer generating floor plans from the
specification failed \factDownstreamFailed{} of \factDownstreamTotal{} generations,
because four of the seven bad lines were forms the parser rejects. The grammar's cost
is harder to bound and worse in kind: a constrained decoder sampling under an
over-permissive grammar produces output that is guaranteed well-formed by the grammar
and rejected by the parser, which is the precise inversion of the artifact's purpose.

\subsection{The dividing line is an executable counterpart}

Non-reflexive validation of a machine-consumed description is neither rare nor new. It
is mature wherever the description has something to be run against: Pact executes a
generated contract against a live provider~\cite{pact_contract_testing}, Schemathesis
executes a schema against a running server~\cite{schemathesis}, protobuf has a
conformance suite~\cite{protobuf_conformance}, and Kubernetes fuzzes real values through
its generated conversion code so that a conversion which drops data
fails~\cite{k8s_roundtrip}. Doctests are the limiting case: a doctest is
executable \emph{by construction}, because it is code.

What these share is not authorship --- Pact's contract is generated, not hand-written
--- but the availability of an executable counterpart. Table~\ref{tab:counterpart}
splits the field on that line.

\begin{table}[t]
\caption{Whether a machine-consumed description has something to run against}
\label{tab:counterpart}
\centering\footnotesize
\begin{tabular}{@{}p{0.44\columnwidth}p{0.44\columnwidth}@{}}
\toprule
\textbf{Has an executable counterpart} & \textbf{Has none} \\
\midrule
OpenAPI schema + a live server; a Pact contract + a real provider; a doctest; a
protobuf conformance suite; generated conversion code + a fuzzer &
a constrained-decoding grammar; a specification page; a bundled context file; a
protocol handshake identity; a vendored documentation resource \\
\addlinespace
\emph{validation: mature, mainstream} & \emph{validation: absent} \\
\bottomrule
\end{tabular}
\end{table}

The right-hand column is the LLM-facing artifact class, and it is where we measure.
Kubernetes' round-trip fuzzer is worth dwelling on: we found it while trying to refute
this dividing line, and it supports it. Generated \emph{code} gets fuzzed because it can
be run. A grammar, a specification page and a handshake string cannot be run, and get
nothing.

\section{An Ecosystem Census}\label{sec:census}

Of the \factEcoApplicable{} npm packages on the Model Context Protocol
registry~\cite{mcp2026spec} whose self-reported identity can be settled by reading their
shipped code, \factEcoDriftConfirmed{} --- \factEcoDriftConfirmedPct{} --- state at handshake
time a version that once tracked the package and has since been left behind. Each of those is
corroborated case by case against the package's own release history by the procedure in
Section~\ref{sec:corrobstep}. Counting instead every literal that merely disagrees with the
package version, whether or not it could be corroborated, gives \factEcoDriftRaw{} ---
\factEcoDriftRawPct{}. We report that second figure as an upper bound and not as the result.

\subsection{The frame, and why this artifact}

An MCP server is a package whose purpose includes describing itself to a language
model~\cite{anthropic2024mcp}. It states a name and a version to its host at handshake time,
and ships tool descriptions, schemas and documentation that a model reads in order to use it.
The handshake version is the smallest member of the artifact class of
Table~\ref{tab:counterpart}: one string, consumed by a machine, with nothing to be executed
against. It is also the member of that class we found most tractable to measure at ecosystem
scale, because a registry enumerates the population and every candidate ships its answer in
its own tarball.

The frame is every server listed on the official registry that distributes an npm package.
Metric~1 below is a \emph{census} of all \factEcoPackages{} unique packages in that frame
rather than a sample of it, so it carries no sampling error; Metric~2 is a sample and is
reported separately in Section~\ref{sec:metrictwo}, never merged into Metric~1's rate. The
harness has no dependencies, the per-package dataset is retained, and a separate script
recomputes every figure quoted here from that dataset and asserts eight internal identities.

\subsection{Denominators}

Table~\ref{tab:ecofunnel} gives the population at each stage. Three attritions between the
enumerated registry and the applicable set are worth naming, because all three bias the result
downward. Some registry entries are remote endpoints or packages published elsewhere and
declare no npm package at all. Of the tarballs that were read, a small number ship no
JavaScript or TypeScript to scan, and a much larger group ships code in which no handshake site
could be located --- overwhelmingly minified bundles where the SDK constructor has been renamed
and the argument object restructured past recognition. Finally, packages whose version
expression could not be settled statically, because it is supplied by a caller or a
configuration object, are excluded rather than guessed in either direction.

\begin{table}[t]
\caption{Metric 1: each row is the population the next row filters}
\label{tab:ecofunnel}
\centering\footnotesize
\begin{tabular}{@{}lr@{}}
\toprule
\textbf{Stage} & \textbf{Count} \\
\midrule
Server entries enumerated (latest of each) & \factEcoServers{} \\
Unique npm packages declared among them & \factEcoPackages{} \\
\quad version \emph{derived} at runtime --- cannot drift & \factEcoDerived{} \\
\quad version a hardcoded \emph{literal} --- can drift & \factEcoRetyped{} \\
\textbf{Applicable} (derived $+$ literal) & \textbf{\factEcoApplicable{}} \\
\midrule
Disagrees with the package version (upper bound) & \factEcoDriftRaw{} \\
\textbf{Corroborated as left behind} (headline) & \textbf{\factEcoDriftConfirmed{}} \\
\bottomrule
\end{tabular}
\end{table}

\subsection{Derived versus retyped, measured}

The split between those two middle rows is the field measurement of
Section~\ref{sec:mechanism}'s precondition, and it is the single most important number in this
paper. \textbf{Of the packages whose self-reported identity can be resolved statically,
\factEcoDerivedPct{} derive it and \factEcoRetypedPct{} retype it.} A derived version is read
from the package's own manifest at load time and cannot be wrong about the package; a retyped
one is a literal a person wrote next to a constructor, and it is wrong as soon as the next
release ships. The reflexive-gate ecosystem of Section~\ref{sec:intro} is built for the first
case, its soundness depends on being in the first case, and in this population the first case
is the minority by an order of magnitude.

We note what the split does \emph{not} say. It is not a claim that
\factEcoRetypedPct{} of these packages are wrong today --- most of the retyped ones are, but
that is the separate measurement above. It is a claim about which failure is available: for
\factEcoDerived{} packages the only reachable defect is staleness of the manifest itself, and
for \factEcoRetyped{} the description and the subject are free to disagree indefinitely.

\subsection{Corroboration}\label{sec:corrobstep}

Observing that a literal happens to be some earlier version of its package is not sufficient
evidence of drift. In a package with hundreds of releases any short version string appears
somewhere in its history by coincidence, and some projects deliberately version an embedded MCP
surface independently of the package that carries it.

So for each candidate the harness downloads the release the literal \emph{names} and asks the
decisive question: at that version, did the handshake say exactly this? If it did, the literal
once tracked the package version and has since been frozen, and the drift is real rather than
an independent scheme. If it did not --- if that release's handshake said something else, or
had no handshake, or no literal --- the case is excluded.

\factEcoIndependentVersioning{} cases were excluded on precisely those grounds: at the version
their literal names, the handshake said something different, so the literal is evidence of a
separate versioning scheme rather than of neglect. Those exclusions are the check working. They
are also the reason a hardcoded handshake version cannot be assumed wrong on sight, and the
reason we lead with the corroborated figure: the difference between \factEcoDriftRawPct{} and
\factEcoDriftConfirmedPct{} is not noise but the cost of insisting on per-case evidence.

\subsection{A category that cannot be corroborated}

\factEcoNeverPublished{} packages state a well-formed version that has never been any published
version of that package. \factEcoNeverPublishedScaffoldPct{} of these values are a common
scaffold default, which is consistent with a placeholder that was never revisited --- though we
did not inspect commit histories and so cannot say when any of them was written. These are not
values that were once true and went stale; they were never true. There is no
earlier tarball to check them against, so they are \emph{not corroborable by construction} and
are excluded from the headline, even though a value that was never correct is arguably the
worse failure. The headline therefore understates.

Two smaller categories are worth recording because they are not staleness either. A handful of
literals are not versions at all, including one package that ships its own uninterpolated build
template as its identity: a substitution that never ran, published to the registry and served
to every host that installs it. In a further handful the literal names a version \emph{newer}
than the package --- a bump applied to the description and not to the release.

Separately, \factEcoInternalDisagreement{} packages contain two or more different hardcoded
version literals in their own shipped code, so their files disagree with each other about what
version they are, independently of whether any of them matches the manifest. Which answer a
host receives depends on which entry point it launched.

\subsection{Distribution, not just the median}

The median corroborated case is \factEcoMedianReleases{} releases and one day behind. On its
own that is unremarkable: a maintainer bumped the manifest, published the same day, and left
the literal. Quoting only the median would understate the phenomenon and quoting only the tail
would overstate it, so both: \factEcoGeFive{} packages are five or more releases behind their
own stated identity, \factEcoGeTen{} are ten or more, and the confirmed extreme is
\factEcoMaxReleases{} releases spanning \factEcoMaxDays{} days. Roughly half the corroborated
population is a same-day slip and roughly half is a package genuinely out of date about itself.

\subsection{False positives}

An automated count of \factEcoDriftConfirmed{} cases that nobody has read is the kind of number
this paper argues against, so the count was checked in two ways.

First, over the whole confirmed population: every positive must carry its own evidence, meaning
the recorded source excerpt contains the claimed version verbatim or the record names the
constant it resolved from. All of them do, and the verification script fails if that ever stops
holding. This is a completeness property of the record rather than of the judgement, but it
establishes that no confirmed positive rests on an unrecorded inference.

Second, a deterministic sample of \factEcoFpSample{} confirmed positives was read as source by
a reviewer working independently of the harness, yielding \factEcoFpFound{} false positives.
The rule of three at that sample size with no observed events gives
\factEcoFpUpperBound{} as an upper bound at the conventional confidence level. That bound is loose, and we state it as the honest limit of a hand
review at this size rather than as a precise rate. Errors of the detector concentrate in the
never-published stratum and in the minified-bundle fallback, neither of which contributes to
the headline.

\subsection{One package that is the whole argument}

\texttt{@shiplightai/mcp} ships a single bundle that contains both of its self-descriptions. One
literal is the version reported by its \texttt{-{}-version} flag, and it is correct. The other
is the version passed to the MCP handshake, and it is wrong. The two sit in one file, written at
the same time, describing the same artifact at the same instant.

\textbf{The human-facing self-description is right and the model-facing one is wrong, and
nothing distinguishes them except who reads them.} Nobody filed a bug about the handshake
because no person sees it; the flag is correct because people run it and would notice. That is
the thesis of this paper in one package, and it was not constructed for the purpose --- it is
the current published release.

\subsection{Metric 2: vendored resources}\label{sec:metrictwo}

The second metric asks a different question of a different population, with its own
denominators, and its rate must never be combined with Metric~1's. An MCP server may vendor
documentation, schemas or grammars into its tarball at pack time; those copies are then frozen
while the project's own sources move. To compare them, a package must declare a public
repository, that repository must be reachable, and it must carry a tag matching the published
version.

That last requirement is the binding one. The sample was drawn from packages declaring a
parseable repository, and fewer than half of the drawn packages could be compared at all:
repositories are private or deleted, and many projects publish releases without tagging them.
\factEcoMTwoCompared{} packages survived to a comparison. Of those,
\factEcoMTwoDiffer{} --- \factEcoMTwoDifferPct{} --- ship at least one vendored text file that
differs from the same file in the project's own repository at the matching tag.

We report \factEcoMTwoDifferPct{} as a raw upper bound. It counts any textual difference,
including formatting and build-time rewriting, and hand verification of the differing cases
reduces it substantially. The attrition and that ambiguity together are why Metric~2 is
supporting evidence for the same mechanism in a second artifact and not a headline.

\subsection{What the census cannot see}

Every bias we can identify points downward. Packages with no locatable handshake are dropped
rather than assumed to drift, and they are disproportionately minified bundles, which is
neither a random nor an obviously cleaner subpopulation. Unresolvable version expressions are
dropped. A package counts as a positive only if none of its literals matches its own version.
The never-published category, which is the more serious failure, is excluded from the headline
entirely. And a package whose handshake version happens to be correct today may be retyped and
one release from being wrong; the census records it as clean.

\section{Four Independent Codebases}\label{sec:corroboration}

The census measures one artifact at scale. It does not show what the same failure looks like
inside a compiler, so we examined four third-party floor-plan compilers written by other
people: two independent implementations of one language (a TypeScript original and a Rust
port by a different author), an IFC toolkit, and a Typst drawing library. Each was cloned at a
pinned commit, and every finding is reproduced by a committed script that names the specific
evidence it depends on. Table~\ref{tab:corrob} lists the six instances.

We sort them into three sub-classes. \textbf{A}: a system's description of itself disagrees
with its implementation. \textbf{B}: a declared diagnostic that can never fire, or a check that
computes a value and discards it. \textbf{C}: a guard or a shipped surface that is never
executed, so nothing about it is verified. Only the first is drift in the sense of
Section~\ref{sec:mechanism}; we keep the other two separate rather than folding them into one
headline count.

\begin{table*}[t]
\caption{Six instances across four third-party compilers, each at a pinned commit}
\label{tab:corrob}
\centering\footnotesize
\begin{tabular}{@{}llp{0.42\textwidth}ll@{}}
\toprule
\textbf{Project} & \textbf{Sub-class} & \textbf{Finding} & \textbf{Verdict} & \textbf{Evidence} \\
\midrule
PlanScript (TS) & B &
\texttt{isPolygonClosed} computes a first point, a last point and a tolerance, then returns
true unconditionally; it is never called, and the error code it exists to raise has no raise
site anywhere & confirmed, stronger than seeded & reading \\
\addlinespace
PlanScript (Rust) & B &
three declared error-code variants are never constructed anywhere in the crate or its tests,
while a design document advertises two of them as typical compiler errors & confirmed, broader
than seeded & execution \\
\addlinespace
PlanScript (Rust) & A, C &
placement order is taken from a hash map, so the solver is non-deterministic; determinism is
promised in four documents and tested in none & confirmed & execution \\
\addlinespace
PlanScript (both) & C &
the connectivity check counts a shared wall as a connection, and the one test written to prove
that doors establish connectivity passes with the door branch deleted & partially confirmed;
drift framing refuted & execution \\
\addlinespace
ifc-lite & C &
the wall clock is written into every exported drawing's title block, while the file's only
byte-identity assertion runs on the path that excludes the title block & confirmed & execution \\
\addlinespace
arch-plotter & C &
a mistyped anchor becomes a zero-length move and the wall silently disappears; the library
raises no diagnostic of any kind, and its two test files have never been executed & confirmed
& execution \\
\bottomrule
\end{tabular}
\end{table*}

Four of the six are backed by a run rather than by reading the source. That was deliberate.
This paper argues that reading a surface is not verifying it, so evidence gathered by reading
would be gathered by the method the paper criticises. Reading tells you that a hash map decides
placement order; running the solver thirty times on one input and getting thirty distinct
buildings tells you it matters. Reading tells you a date is formatted from the system clock;
exporting the same model under two stubbed clocks and diffing the bytes tells you the drawing
is not reproducible, and that a locale-dependent format makes it vary between two colleagues
exporting at the same instant. In two of the four cases the executed result is strictly
stronger than what the source supported, and in one it corrected the claim we started with.

Each script was also negative-controlled: a scratch copy of the clone was patched to remove the
evidence, and the script re-run to confirm it then fails with a message naming what went
missing. All six pass on the real clone and fail on the patched copy. A reproduction script
that has never been run is the defect class this paper documents, and a script that cannot fail
is that defect one level up.

\subsection{A defect that survived a reimplementation}

The strongest single result is the pair of PlanScript ports. \texttt{E101} --- the code for an
unclosed polygon --- is dead in both. In the TypeScript original the function that would raise
it returns true unconditionally and is never called; in the Rust port, written later by a
different author, the enum variant exists and is never constructed. The port re-typed the error
catalogue faithfully, including the entries with nothing behind them. \textbf{The defect
survived a cross-language reimplementation by a different person}, which makes it a property of
maintaining a catalogue by hand rather than of one author's carelessness.

The two catalogues then diverge in a way that is more informative than the agreement. The
TypeScript version has four dead codes and the Rust port has three: the port revived one code
that is dead upstream, and inherited the other three. Overlapping-but-different sets are what
one would predict if each catalogue were hand-maintained rather than derived from raise sites,
because two people transcribing the same list make different omissions. A derived catalogue
cannot produce that pattern at all, so the divergence discriminates between the two
explanations rather than merely illustrating one of them. None of the four
dead codes is named by any test in either repository, which is consistent with codes that cannot
fire and is the reason no gate reported them.

\subsection{Refutations}

Four things we looked for were not there, and a section that reports only confirmations is not
evidence.

\emph{The drift framing of the connectivity check is wrong.} We expected a description-versus-%
implementation disagreement in \texttt{rooms\_connected} and did not find one: the language
reference documents the shared-edge rule accurately, and both implementations match their own
documentation. What is actually wrong is narrower --- the assertion's name and its error code
promise circulation connectivity while the check measures geometric adjacency --- and it belongs
in sub-class C rather than A. The instance stays in the table with its framing corrected rather
than being dropped. The vacuous test underneath it is a genuine finding and was found only
because the original claim was checked instead of accepted.

\emph{The dead-code pattern is not universal.} A sweep of ifc-lite's MCP error catalogue found
all eleven declared codes in active use, between one and forty-six call sites each. A project
can evidently keep a small catalogue honest, and the claim is therefore falsifiable and was not
rubber-stamped.

\emph{Not every generated artifact can drift.} PlanScript's parser is generated by a parser
generator from a committed grammar file, and regenerated on every test run and every build; the
generated parser is not committed at all. There is nothing to go stale. This is the derived case
of Section~\ref{sec:mechanism} implemented properly, in one of the same repositories that
carries three dead error codes.

\emph{A shell pattern we expected was absent.} The example-building scripts in both PlanScript
repositories use \texttt{set -e} correctly, so the ``script that cannot fail'' pattern is not
present there. What we found instead in the Rust port is different in kind: two of its five
shipped example inputs no longer solve, so the script that regenerates the examples cannot
complete, and the committed outputs for those two are artifacts of a build that no longer runs.
That project's non-determinism also forecloses the reflexive gate entirely --- a
regenerate-and-compare step is impossible where regeneration produces a different building each
time --- which is worth stating because it is the one case where the gate's absence is not an
omission the project could simply correct.

\subsection{What is conceded on this axis}

Requiring that a tool's declared diagnostics can actually fire is not our idea and is already
shipped and CI-enforced at the top of the field: rustc's \texttt{tidy} check demands that every
error code be documented, be triggered by a test, and be emitted by the
compiler~\cite{rustc_tidy_errorcodes}. The mechanism these instances need already exists in a
production compiler.

Our contribution on this axis is therefore the \emph{distribution}, not the idea. The practice
exists where somebody paid for it and is absent below that tier: four dead codes in one
compiler, three in its port, a design document advertising two of them to readers, and no test
in either repository naming any of them. What separates rustc from these projects is not
knowledge of the technique but whether anyone made the catalogue answer to the code.

\section{One System in Depth}\label{sec:depth}

The census counts one string per package and the third-party instances are one defect each.
Neither shows what an entire self-describing surface looks like when it is checked for the
first time. The subject system is a production floor-plan compiler with a command-line
interface aimed at language models: it ships a specification page, a constrained-decoding
grammar, JSON schemas, a bundled context file, and a protocol shim that vendors copies of all
of them. Every one of those artifacts had a reflexive drift gate. Section~\ref{sec:twogenerators}
reports what happened when two of them were first executed against the parser instead of
against themselves; this section reports the measurements that followed, including two that
were themselves defective.

\subsection{An agreement corpus for the grammar}

The grammar is checked against the parser by a corpus of complete plan sources, each run through
both. The corpus is \factGbnfCorpus{} sources chosen to sit on clause-ordering and arity
boundaries, of which \factGbnfParseable{} are accepted by the parser and \factGbnfRejected{} are
rejected. That split is itself asserted, so the corpus cannot be satisfied by a grammar that
accepts everything or one that rejects everything.

For each source the property is a biconditional: the grammar derives it if and only if the
compiler parses it. Against the last grammar that was ever published, agreement is
\factGbnfAgreeBefore{} of \factGbnfCorpus{}. Against the current grammar it is
\factGbnfAgreeNow{} of \factGbnfCorpus{}. The \factGbnfDisagree{} disagreements are not
symmetric: \factGbnfOverDerive{} are over-derivations, where the grammar admits a form the
parser rejects, and \factGbnfUnderDerive{} are under-derivations, where the parser accepts a
form the grammar cannot produce. The over-derivations are the serious direction, because a
constrained decoder sampling under that grammar emits them as well-formed by construction and
the parser then refuses them.

The methodological point is not the ratio but where the expected value comes from.
\textbf{The expected verdict for each source is computed by calling the compiler, on every run,
so there is no expected column to edit.} A conventional table of source-and-expected-result
pairs can be made green by changing the expectation, and a maintainer under time pressure
usually can; here the only two ways to make a red case green are to fix the grammar or to
change the language. This is the difference between a test that records a belief about the
system and one that asks the system.

\subsection{Executing the specification}

The same treatment applied to the prose specification is a suite that extracts every form the
document teaches and every illegality it declares, then runs them. As shipped,
\factSpecFormsCompile{} documented forms must compile and \factSpecFormsRefused{} documented
illegalities must be refused \emph{with the catalogued code the document names for them} ---
refusal alone is not enough, because a form can be rejected for the wrong reason and still look
correct from outside. The suite is bound to the compiler's keyword set, so a keyword with no
documented form fails it rather than being silently undocumented.

Executable documentation is not new and we make no claim on it: Rust, Go and Python all run
their documentation examples by default~\cite{rust_doctests,go_testable_examples,python_doctest},
\texttt{@tComment} turned Javadoc into runnable assertions in
2012~\cite{tan2012tcomment}, and CASCADE generates tests directly from
documentation~\cite{kiecker2026cascade}. What is worth reporting is where it was absent. The
document these forms come from is read by no human in the normal course of events, so it had
none of the pressures that produce doctests --- no reader to be confused, no example to be
copied into a tutorial --- and it carried its errors for three releases behind a green gate.

\subsection{A mutation score for two modules}

Two modules of the subject system had no direct caller in the test suite: the arithmetic behind
the measured-deficit diagnostics, and the exact-isometry layer that carries a placed component's
geometry into plan coordinates. Both were reached only incidentally, through whole-plan tests.
To measure what those tests were actually checking we ran a mutation experiment in the standard
sense~\cite{demillo1978hints,jia2011mutationsurvey}: \factMutants{} hand-authored mutants,
\factMutantsMeasure{} in the first module and \factMutantsFrame{} in the second, each a
find-and-replace pair whose target must match exactly once in its file, asserted before anything
is written.

Each mutant is applied, the module's dedicated test file is run, and if that file still passes,
a fixed set of whole-plan suites that reach the module is run as well. \factMutantsKilled{} of
\factMutants{} mutants are killed, a mutation score of \factMutantScore{}. The
\factMutantsSurvived{} survivors are all judged equivalent, with the argument stated per
survivor: two unconditionally, because the mutated comparison cannot be reached with a
distinguishing value, and three under an invariant that their only caller establishes. Equivalent
mutants are a known and undecidable-in-general obstacle to interpreting a mutation
score~\cite{papadakis2019mutationadvances}, so we state the reasoning rather than assert the
classification, and we note the weaker form it takes: a survivor is equivalent \emph{under its
caller's invariant}, which is a property of the current program and not of the function.

Two facts about the run are worth more than the score. The whole-plan suites killed no survivor
that the dedicated file did not already kill, so widening the suite set does not move the number
--- the incidental coverage that made these modules look tested compensates for nothing. And an
earlier version of this experiment surfaced two survivors that were \emph{not} equivalent: both
in the reflection path of the transform layer, each demonstrated by a plan whose rendered bytes
the mutation changes while the suite stays green. Those are now killed, which is why the
survivors reported here are all equivalent.

\subsection{A worked example: the run that measured nothing}

An earlier run of this same harness over the same subject reported \factMutantsVoidKilled{} of
\factMutants{} --- \factMutantScoreVoid{} --- and it is void.

The harness works by writing a wrong version of a source file into the tree, running the tests,
and putting the file back. While it runs, the repository contains a compiler that is wrong on
purpose, and a mutated file is indistinguishable on disk from an ordinary uncommitted edit.
Partway through that run, a third party noticed a source file that looked corrupted and restored
it from the last commit. The interference was the session lead directing this work --- the person
who commissioned the measurement, acting reasonably on what the file system showed. Naming that
is the point: no rule about who may touch the tree would have prevented it, because the
information needed to distinguish a live measurement from a leftover was not available to the
person who had to decide.

Two mutants were then counted as survivors purely because their mutation had been removed from
underneath them. The resulting score is plausible: it sits slightly below the true one, in the
direction and by the sort of margin a reader would expect a real result to sit, and there is
nothing in the number to suggest anything went wrong.

The artifact of that run records \texttt{treeCleanAfterRun: true}. The check was honest and it
was not lying: the tree was clean when the run ended, because the interference had restored the
file. It answered the question it was asked. It was not asked whether the file had been clean
\emph{throughout}, which is the question that mattered, and no reader of the artifact could tell
the difference. \textbf{A measurement corrupted by concurrent interference produced a
plausible-but-wrong number and self-certified as clean.} We report this as a worked example
rather than as a result, and the earlier score appears nowhere in this paper as a measurement of
anything.

The remedy is the same shape as the rest of the paper. The harness now records a marker before
it does anything at all, carrying which phase it is in, and answers ``is a run executing?''
rather than ``is the tree mutated?'' --- two questions that look identical on disk. The
instruction is to ask the tool, not to inspect the file system. An earlier version of that
instruction told readers to check for the marker file's existence, and there was a window in
which a live run had not yet written it; a reader following the instruction during that window
would have concluded that a live measurement was a leftover.

\subsection{Surfaces that shipped without ever being executed}

The last group is not about descriptions but about the same blindness applied to behaviour, and
it is what the depth study turned up once the executable-check habit was established.

The subject system's file-watch command had not watched for \factWatchBrokenReleases{} releases,
across which the machine-readable command manifest continued to advertise that it recompiles on
save. A refactor had wrapped the command so that the process exited immediately after arming the
watcher. Nothing was wrong with the manifest's generator, the manifest was regenerated and
compared on every build, and the command had no end-to-end test of any kind: the gate guarding
the description of the feature was green throughout, because the description was never the thing
that was broken.

The system's flagship determinism property is a second instance, and it is a generator-coverage
failure of the kind the property-testing literature already
names~\cite{groce2012swarm,goldstein2021judgecover}. The property fed a general string generator
into the body of a plan and asserted that compiling twice produced identical bytes. Across
\factOldPropSamples{} samples it produced \factOldPropGeometry{} walls, rooms, openings or
fixtures: essentially every sample was a parse error, and the property was asserting that two
compilations of an invalid plan agree. It passed continuously and its subject was empty.

The replacement generates plans by construction, drawing every closed value set from the module
that owns it, and renders geometry in every one of its \factNewPropSamples{} samples. The
difference is measurable rather than asserted: an iteration-order bug planted in the wall-lowering
pass is killed by the new property within \factPlantKilledRuns{} runs, and survived
\factPlantSurvivedRuns{} runs of the old one. That comparison is the vacuity-detection procedure
of Beer et al.~\cite{beer1997vacuity} --- perturb the subject and require the guard to
notice --- applied to a property test, and it is the only evidence we accept for a guard in this
paper.

\subsection{A fifth self-description, found while writing this}\label{sec:fifth}

The four artifacts above were checked in one pass and are a closed set. A fifth surfaced during
the writing of this paper, in the same repository, and is reported because it was not selected
retrospectively --- and because it inverts the expected direction.

A documentation site rendered the language's own source in one flat colour, which is nearly
everywhere it appears: every unmarked code fence on that site becomes a live editor. The cause was
a hand-typed keyword set of \factHandTypedWords{} words inside a component, carrying a comment
asking the next person to keep it in step with the compiler's table. This is the retyped generator
of Section~\ref{sec:mechanism} without even a gate over it, and its failure mode is the one that
makes this class hard: a highlighter that omits a word does not fail --- it draws that word as an
identifier, which reads as a typo rather than as a defect.

The expected repair is to derive it, and that is what was done. **What was not expected is that the
hand-typed copy was more correct than the source it was supposed to track.** Two words, \texttt{street}
and \texttt{hemisphere}, were present in the hand copy and absent from the compiler's keyword
table, where they had been missing since \factUntabledSince{} --- accepted by the parser and
documented in the specification's own grammar line for that construct, and therefore drawn as bare
identifiers by every renderer for \factUntabledReleases{} releases. Deriving the artifact would
have \emph{regressed} those two words. What found them was the guard shipped alongside the
derivation: an assertion that every member of every closed value set also appears in the keyword
table. The table went from \factKeywordsBefore{} entries to \factKeywordsAfter{}.

The lesson refines Section~\ref{sec:remedy}'s remedy rather than confirming it. Deriving a
description does not make it correct; it makes its correctness \emph{equal to the source's}. That
is a gain only because the act of deriving forces the source's incompleteness into view --- here,
a hand copy scoped by its own comment to ``everything one example file uses'' had been quietly
carrying two facts the source of truth did not.

\section{Remedy: Provable Non-Vacuity}\label{sec:remedy}

The remedy for a retyped generator is to derive instead --- which is
DRY~\cite{hunt1999pragmatic}, and ours only in the sense that we arrived at it after
violating it. What we can offer is narrower and, we think, more useful: the reflexive
gate's precondition can be made checkable, and the guards that replace it can be
required to prove they are not vacuous.

\subsection{Four mechanisms}

\textbf{Single-source the closed vocabularies.} Every finite set of values the language
admits --- element kinds, enum words, door kinds, paper sizes --- lives once, and every
description of the language interpolates from it. The subject system carries eight such
sets and an \texttt{assertVocabRendered} guard that throws when a member of a set has no
rendering in the artifact being generated. This converts the generator from retyped to
derived for exactly the facts most likely to grow, and it fails loudly when a new value
is added without a home.

Its second effect is the one we did not anticipate, and it is why this mechanism is listed
first. Deriving a description does not make it correct --- it makes its correctness equal to
the source's, which is an improvement only if the source is complete. What the guard actually
buys is that incompleteness stops being invisible: in Section~\ref{sec:fifth} it is what
revealed that a hand-typed copy had been carrying two values the source of truth lacked, so
that deriving the copy would have lost them. A single-source rule with no completeness check
moves the defect rather than removing it.

\textbf{Compute the expected value rather than write it down.} The GBNF agreement corpus
of Section~\ref{sec:depth} pairs each of \factGbnfCorpus{} plans with a verdict taken
from the compiler \emph{on every run}. There is no expected column in the fixture, so
there is nothing a maintainer can edit to make a red suite green; the corpus can only be
satisfied by fixing the grammar or by changing the language. This is the structural
difference between a test and a snapshot, and it matters most exactly where a snapshot
is most tempting.

\textbf{Execute the documentation.} The reference is a test corpus:
\factSpecFormsCompile{} documented forms must compile and \factSpecFormsRefused{}
documented illegalities must be refused \emph{with their catalogued code}. This is
\texttt{@tComment}'s idea~\cite{tan2012tcomment} applied to a generated artifact, and
rustc's \texttt{tidy} check~\cite{rustc_tidy_errorcodes} is the strongest existing
instance of it; we claim the application, not the technique.

\textbf{Require every guard to prove it is not vacuous.} A guard that has never been
observed to fail is a guard nobody has reason to believe. The available proofs are a
planted fault the guard must reject, an anti-vacuity assertion that a fixture really is
faulty before it is healed, and --- for a corpus --- a non-vacuity case asserting that
the corpus contains both accepted and rejected members, so it cannot be satisfied by a
recogniser that says yes to everything.

\subsection{The invariant we converged on was already universal}

The single-source rule is not an invention. It is the property the derived-artifact
ecosystem relies on everywhere: API Extractor reads a symbol table, Gazelle reads real
imports, \texttt{protoc} reads the schema. The finding is not that the rule is good ---
it is 27 years old --- but that it is \emph{universal in practice, stated nowhere, and
unenforced by the very gate meant to protect the artifacts that depend on it}. A project
can violate it and keep a green build indefinitely, which is what the subject system did
for three releases and what \factEcoRetypedPct{} of the packages in
Section~\ref{sec:census} are doing now.

\subsection{What the precondition check would look like}

We do not have a general detector, and we do not claim one. But the question a reflexive
gate should be able to answer about itself is small and concrete: \emph{does this
generator's output depend on any fact that a human typed into the generator?} For the
narrow case of closed value sets the answer is mechanisable, and the
\texttt{assertVocabRendered} guard above is one implementation. For prose that states a
grammar rule in English, it is not; there, executing the documentation is the only
remedy we know.

\section{This Paper's Own Toolchain}\label{sec:self}

The strongest evidence we can offer that this failure mode is ordinary rather than a
story about one careless project is that the tooling built to support \emph{this paper}
produced seven vacuous guards in the course of writing it. They are listed in the order
they occurred, with who or what caught each.

\begin{enumerate}
\item A bibliography set-comparison whose pattern collapsed and matched nothing,
  reporting every entry as missing. Caught only because a different check contradicted
  it. Its author's own note is the sharpest sentence any of us wrote: \emph{``I ignored
  it because grep contradicted it --- but had it reported success I would have believed
  it.''}
\item A page-limit checker that scraped the PDF for a marker absent from compressed
  object streams, fell through to a fallback that matched an unrelated number, and
  reported a one-page paper as 47 pages. It would have passed any limit. Caught because
  the number was absurd, not because anything failed.
\item A fault injector whose planted citation lost a backslash in transit, so the
  non-vacuity proof it was performing proved nothing while reporting success. Caught by
  checking the bytes on disk rather than the tool's own report.
\item A stale \texttt{.tex} silently recompiled against a bibliography that had grown
  beneath it. \textbf{Caught by a guard} --- the only one of the seven that was.
\item A superseded paragraph left standing directly above its own correction in the
  planning document, with nothing marking which was which. Caught by a reader.
\item A mutation harness that left a planted fault in the source tree while its result
  file recorded \texttt{treeCleanAfterRun: true}. Caught by routine inspection. Its
  replacement's safety marker was then documented in a \texttt{SAFETY.md} that would
  have caused the corruption it was written to prevent, because the marker was written
  after the first mutation rather than before the run; following that document would
  have destroyed a live measurement for the second time.
\item A citation checker that short-circuited on \texttt{\textbackslash nocite\{*\}} ---
  correct in principle, since that macro defines every key, and wrong in practice,
  because a \emph{misspelled} key is undefined even then. It waved one through while the
  build reported success. This was the guard written to catch \#4.
\item A banned-claim checker, written to enforce the retractions of
  Section~\ref{sec:concede}, that matched multi-word phrases against unwrapped source.
  \LaTeX{} wraps at the margin, so a banned phrase split across a newline evaded it. Three
  phrases were planted; it caught the two that happened to sit on one line and missed the
  third. \textbf{It was real and partial, which is worse than absent, because partial
  coverage reads as coverage.}
\end{enumerate}

Two of eight were caught by a gate; the rest were caught by a number looking wrong, by
checking bytes instead of a report, by a reader, and by routine inspection. Every one of
the eight was found because a fault was deliberately introduced and the guard was asked to
notice --- which is Section~\ref{sec:remedy}'s requirement, applied to the guards written
for this paper, and the only reason the list is eight long rather than shorter and wrong.

One contrasting case is worth reporting because it is the design working. The page-limit
checker could not parse the label reference emitted by a different document class, whose
counter is wrapped in a macro the checker's pattern could not span. It did not pass. It
reported that it could not verify what it claimed to check, and exited non-zero. A guard
that cannot see its subject and says so is not a member of this list; it is the behaviour
the other eight lacked.

There is also an observer effect worth a sentence. Committing the input behind one of this
project's transcripts required a note recording its provenance; adding that note to the
input file moved the byte offsets the transcript reported, silently invalidating it. The
note explaining an artifact changed the artifact it described, and the provenance now lives
beside the file rather than inside it.

We report all of this not as self-deprecation but as the paper's most direct evidence.
These guards were written by people actively arguing that this failure is ordinary, while
they were arguing it, in a repository whose continuous integration was green throughout.

\section{Threats to Validity}\label{sec:threats}

\textbf{Construct.} We measure whether a machine-consumed self-description is
\emph{true}, not whether being false caused harm in any particular deployment. For the
handshake version the harm is plausible but unquantified: a host that logs or gates on a
server's reported version is misled, and we did not measure how many do. For the grammar
and specification the downstream cost was measured once, in one consumer
(Section~\ref{sec:twogenerators}), and one measurement is not a rate.

\textbf{Internal.} The census depends on statically resolving a value out of published
JavaScript. Packages where that resolution fails are excluded rather than guessed, and
the exclusion is not random: reading a version from \texttt{package.json} is itself
dynamic, so if dynamic resolution correlates with better engineering, the true retyping
rate across \emph{all} packages is lower than the rate among applicable ones. We
therefore quote the rate among applicable packages and say so.

\textbf{External.} Metric~1 is a census of registry-listed servers with an npm package,
so it carries no sampling error --- but that frame is not ``all MCP servers'': servers
distributed via PyPI, Go, OCI or as remote endpoints are absent, as are unregistered
ones. Metric~2 \emph{is} a sample and its rate carries sampling error accordingly. The
four corroborating codebases are all floor-plan tooling, chosen because the subject
system is; whether the same rates hold in other domains is untested. One counterexample
(ifc-lite's error catalogue, entirely live) shows the pattern is not universal.

\textbf{Conclusion.} The mutation score is over a hand-authored mutant set, which
encodes its author's idea of what can go wrong; it is a measurement of that suite
against those mutants and is not comparable to any other project's score, nor to an
earlier unreproducible figure from this project which we deliberately do not cite. The
subject system is single-maintainer and young, and its eval rates are single runs at
$n = \factEvalBriefs{}$ which we report only for orientation and never as a benchmark.

\textbf{Search.} Our negative claims are of the form ``we found no prior work naming
this''. Two holes are worth stating: a major bibliographic database was unreachable
throughout the search, and unauthenticated code search over public repositories was not
usable at scale, so the engineering-practice sweep rests on documentation and a bounded
set of repositories rather than on a corpus study. That sweep already retracted two of
our claims once, which is evidence about the method's yield in both directions.

\section{On Language-Model Co-Authorship}\label{sec:ai}

The subject system was substantially co-authored with language-model agents, and so was
this paper's tooling. We disclose this in accordance with the ACM Policy on Authorship,
and we think it is more than a disclosure.

The hypothesis, stated as a hypothesis: \textbf{a model writes plausible meta-layer
artifacts, and plausibility is precisely what a reflexive gate cannot distinguish from
truth.} A specification paragraph that reads correctly, a grammar production that looks
like the others, a safety document describing a marker that is never written --- each is
locally well-formed, consistent in tone with its neighbours, and wrong. The properties
that make model-written prose good also make it good at passing a check that never
executes it.

What is consistent with the hypothesis: all four of the subject system's
self-descriptions were wrong on first check; the failures were in artifacts nothing
executes rather than in code, which is comparatively well covered; and the seven guards
in Section~\ref{sec:self} were themselves model-written and each looked right.

What would falsify it: a comparable measurement over human-authored projects of similar
age showing the same rates. We have not run it, and the census in
Section~\ref{sec:census} cannot separate the two --- we do not know which of those
\factEcoPackages{} packages were model-assisted. We therefore state the hypothesis and
decline to claim it. The measurement stands without it.

\section{Related Work}\label{sec:related}

Section~\ref{sec:concede} conceded the components; this section places them, and says in each
case what remains.

\textbf{Vacuity.} That a check can pass without its passing depending on what it appears to
check is vacuous satisfaction, formalised for model checking by Beer et
al.~\cite{beer1997vacuity,beer2001vacuity} and Kupferman and
Vardi~\cite{kupferman2003vacuity}, and carried into testing by Ball and
Kupferman~\cite{ball2008vacuitytesting}. Beer et al.\ detect it by mutating a subformula and
requiring the check to notice, which is the planted-fault remedy of
Section~\ref{sec:remedy}, twenty-nine years earlier. We take the concept and the remedy from
this line entirely. What that literature does not supply is which \emph{engineering structures}
reliably produce a vacuous pass; the derived-versus-retyped precondition is an answer to that
question for one widely deployed structure, and it is the only part we claim.

\textbf{Documentation--code inconsistency.} iComment~\cite{tan2007icomment},
\texttt{@tComment}~\cite{tan2012tcomment}, DocRef~\cite{zhong2013docerrors} and Zhou et
al.~\cite{zhou2017directivedefects} detect disagreement between prose and implementation, and
CASCADE~\cite{kiecker2026cascade} now generates tests directly from documentation. The
phenomenon in our Section~\ref{sec:depth} is a member of this family. Two things differ. The
artifacts there are \emph{hand-written} and the remedy proposed is always a detector; ours are
\emph{generated}, a gate already exists over them, and the gate is green. And their consumer is
a reader who can notice; ours is a decoder that cannot.

\textbf{Diagnostic-catalogue completeness.} This is shipped practice rather than literature, and
finding it required looking at build tooling rather than at papers. rustc's \texttt{tidy} check
requires every error code to be documented, triggered by a test, and emitted by the
compiler~\cite{rustc_tidy_errorcodes}; clang ships a script that lists diagnostics defined and
never referenced~\cite{clang_unused_diagnostics}; Roslyn forces every declared analyzer
diagnostic into a release-tracking file~\cite{roslyn_release_tracking}. The nearest research
neighbour is CERTest~\cite{tang2025certest}, a mutation-based approach to compiler diagnostic
testing that explicitly targets missing diagnostics --- but its \emph{missing} is a diagnostic
that should fire for a given program and does not, whereas the instances in
Section~\ref{sec:corroboration} cannot fire for any program at all. The compiler-testing
survey~\cite{chen2020compilertesting} does not carry catalogue coverage as a category. Our
contribution here is a distribution, not an idea: the practice exists where a project could
afford it and is absent below that tier.

\textbf{The same shape in another artifact.} Message catalogues have the problem we describe
and solved it long ago without naming it: the internationalisation ecosystem routinely runs
bidirectional unused-and-missing key detection over translation
files~\cite{i18n_tasks,i18n_unused}. A translation key declared and never referenced is a
declared-but-unreachable entry; a reference with no key is its dual. That this is ordinary
tooling for locale files and absent for diagnostic catalogues is a fact about which artifacts
have been thought worth checking, not about which are checkable. We found no work treating the
two as instances of one problem.

\textbf{Validating a description against its subject.} Pact verifies a generated contract by
executing it against a real provider~\cite{pact_contract_testing}; Schemathesis executes a
schema against a live server and reports divergence without deciding which side is
wrong~\cite{schemathesis}; protobuf ships a conformance suite~\cite{protobuf_conformance};
Kubernetes fuzzes values through its generated conversion code~\cite{k8s_roundtrip}; and Rust,
Go and Python execute documentation by
default~\cite{rust_doctests,go_testable_examples,python_doctest}. This is a mature practice and
we claim none of it. Table~\ref{tab:counterpart} is the observation we do make: it exists
wherever the description has an executable counterpart, and the LLM-facing class has none.

\textbf{Test adequacy.} That a suite can be green and weak is the subject of a long empirical
literature~\cite{inozemtseva2014coverage,zhang2015assertions}, with mutation testing as the
standard adequacy measure~\cite{demillo1978hints,jia2011mutationsurvey,papadakis2019mutationadvances}
and industrial deployment at scale~\cite{petrovic2018googlemutation}. We report our mutation
result in that vocabulary deliberately. The related finding that a property test can be vacuous
because its \emph{generator} never reaches the interesting inputs is generator coverage, studied
under that name~\cite{groce2012swarm,goldstein2021judgecover}. The test-smell
literature~\cite{vandeursen2001refactoringtest,bavota2012testsmells} is sometimes offered as
covering this; it does not. Assertion Roulette concerns many undocumented assertions rather
than none, and that catalogue's concern is maintainability. A test that passes having asserted
nothing is not one of its smells.

\textbf{Measurement validity.} That changing a metric can move a number with no change in the
subject is stated canonically for machine learning by Schaeffer et
al.~\cite{schaeffer2023mirage}, and the construct-validity framing is Jacobs and
Wallach's~\cite{jacobs2021measurement}. The judge recalibration we report in
Section~\ref{sec:depth} is a small instance of their thesis and is offered as corroboration of
it. What survives as ours is procedural rather than conceptual: non-comparability across a
ruler change was pre-committed and enforced mechanically rather than discovered afterwards.

\textbf{Build systems and reproducibility.} Mokhov et al.~\cite{mokhov2020buildsystems} place
the correctness of generated content outside what a build system guarantees, and the
reproducible-builds programme adopts byte-reproducibility as an integrity
property~\cite{lamb2022reproducible}. Both halves of our observation are therefore already
written down. We join them and report what the junction costs.

\textbf{Language models and domain-specific languages.} Anka~\cite{almazrouei2025anka} argues
that a DSL designed for model generation outperforms a general-purpose language, and a survey
covers generation for low-resource and domain-specific
languages~\cite{joel2024dslsurvey}. Constrained decoding is studied given a
grammar~\cite{park2024gad,wang2023grammarprompting,geng2025jsonschemabench}; the correctness of
the grammar artifact itself is assumed throughout, which is the assumption
Section~\ref{sec:twogenerators} violates.

\textbf{The artifact class is new enough to be under-studied.} Publishing documentation
addressed to models rather than people is recent and growing: \texttt{llms.txt} has been
adopted by more than \factLlmsTxtSites{} sites, and agent-tool servers ship bundled context
files as a matter of course. Treude and Baltes~\cite{treude2026contextrot} are the nearest
existing treatment, repurposing documentation-consistency research for AI configuration
artifacts. Their subject is hand-written and has no generator and no gate over it, which is
precisely the configuration our Section~\ref{sec:mechanism} is about; we read their work as
opening the same territory from the other side.

\section{Conclusion}\label{sec:conclusion}

A gate that regenerates an artifact and fails on a difference proves that a generator is
deterministic. Where the artifact is mechanically derived from its subject, that is
enough, and the industry's near-universal use of such gates is not a mistake. Where the
generator retypes a fact instead of deriving it, the same gate is vacuous --- and it
does not check which case it is in, and it looks the same either way.

The artifact class where this now matters most is the one with no executable counterpart
to fall back on: the specifications, grammars, schemas and identities that systems ship
for language models to read. Across a census of \factEcoPackages{} packages,
\factEcoRetypedPct{} retype the identity they could derive, and
\factEcoDriftConfirmedPct{} state a version they are not. The concepts here are old. The
place they have moved to is new, and it is not being watched.

\label{refsstart}
%% `\nocite{*}` was removed once the body carried real citations. It existed only so
%% a skeleton would compile and prove the bibliography builds; leaving it in would
%% ship a reference list of works the paper never cites. Its removal also re-arms
%% build.mjs's cited-vs-defined check, which that macro short-circuits.
\bibliographystyle{IEEEtran}
\bibliography{../refs}

\end{document}
